\documentclass[border=10pt,tikz]{standalone}
\usepackage{fontspec}
\setmainfont{Apple SD Gothic Neo}
\usepackage{tikz}
\usetikzlibrary{positioning,arrows.meta,decorations.pathreplacing}
% _design.tex — shared TikZ design system for all blog diagrams
% Inputs nothing on its own — meant to be \input{} from each diagram .tex
% AFTER \begin{document}, BEFORE \begin{tikzpicture}.
%
% Provides a unified visual language across GoF / DSA / EMC++ / Beautiful series:
%   - Colors: muted, accessible, color-coded by role
%   - Node styles: consistent sizing, rounded, subtle borders
%   - Edge styles: UML-style inheritance, aggregation, plus dashed alternates
%
% Usage:
%   \documentclass[border=10pt,tikz]{standalone}
%   \usepackage{tikz}
%   \usetikzlibrary{positioning,arrows.meta,shapes,fit,backgrounds}
%   \begin{document}
%   \input{../_design.tex}     % adjust relative path
%   \begin{tikzpicture}[blog]  % apply 'blog' style
%   ...

\definecolor{absfill}{RGB}{220,232,247}   % cool blue
\definecolor{absbord}{RGB}{72,118,182}
\definecolor{confill}{RGB}{249,222,222}   % warm red
\definecolor{conbord}{RGB}{197,93,93}
\definecolor{impfill}{RGB}{223,238,222}   % soft green
\definecolor{impbord}{RGB}{99,154,93}
\definecolor{clifill}{RGB}{251,239,205}   % pale yellow
\definecolor{clibord}{RGB}{200,167,72}
\definecolor{hifill} {RGB}{253,224,200}   % accent orange
\definecolor{hibord} {RGB}{210,135,68}
\definecolor{nodefill}{RGB}{250,250,250}
\definecolor{nodebord}{RGB}{120,120,120}
\definecolor{edgec}  {RGB}{80,80,80}
\definecolor{edgesc} {RGB}{145,145,145}
\definecolor{groupbord}{RGB}{200,200,200}

\tikzset{
  blog/.style={
    font=\sffamily\small,
    every node/.style={font=\sffamily\small}
  },
  % --- node roles ---
  cls/.style={
    rectangle, rounded corners=2pt,
    draw=nodebord, line width=0.4pt,
    minimum height=8mm, minimum width=28mm,
    inner sep=3pt, fill=nodefill, align=center
  },
  abs/.style={cls, draw=absbord, fill=absfill, font=\sffamily\itshape\small},
  con/.style={cls, draw=conbord, fill=confill},
  imp/.style={cls, draw=impbord, fill=impfill},
  cli/.style={cls, draw=clibord, fill=clifill},
  hi/.style={cls, draw=hibord, fill=hifill, font=\sffamily\bfseries\small},
  % --- circle (for trees / graph nodes) ---
  vx/.style={
    circle, draw=absbord, line width=0.4pt,
    minimum size=8mm, fill=absfill, inner sep=0pt
  },
  vxc/.style={vx, draw=conbord, fill=confill},
  vxh/.style={vx, draw=hibord, fill=hifill, font=\sffamily\bfseries},
  % --- decision diamond ---
  q/.style={
    diamond, draw=clibord, line width=0.4pt, fill=clifill,
    inner sep=2pt, aspect=2.4, align=center
  },
  % --- edges ---
  inh/.style={                                  % UML inheritance
    -{Triangle[length=3mm,width=3mm,open]},
    draw=edgec, line width=0.5pt
  },
  agg/.style={                                  % UML aggregation
    {Diamond[length=2.5mm,width=2.5mm,open]}-{Stealth[length=2.2mm]},
    draw=edgec, line width=0.45pt
  },
  uses/.style={                                 % directed reference
    -{Stealth[length=2.2mm]}, draw=edgec, line width=0.5pt
  },
  arr/.style={                                  % plain arrow (graphs/trees)
    -{Stealth[length=2mm]}, draw=edgec, line width=0.45pt
  },
  edge/.style={                                 % undirected (trees)
    draw=edgec, line width=0.45pt
  },
  alt/.style={                                  % dashed: similar/alt
    -{Stealth[length=2mm]}, dashed, draw=edgesc, line width=0.45pt
  },
  ret/.style={                                  % dashed creates/returns
    -{Stealth[length=2mm]}, dashed, draw=edgesc, line width=0.45pt
  },
  thr/.style={                                  % threading link (red dashed)
    ->, dashed, draw=conbord, line width=0.5pt
  },
  % --- group / region ---
  group/.style={
    rectangle, rounded corners=4pt,
    draw=groupbord, line width=0.4pt,
    inner sep=10pt, fill=black!2
  },
  glabel/.style={font=\sffamily\bfseries\small, anchor=north west, inner sep=2pt}
}

% Korean font convenience macro — included by every Hangul diagram
\providecommand{\KOR}{}


\begin{document}
\begin{tikzpicture}[blog,
  sda/.style={thick, draw=absbord},
  scl/.style={thick, draw=impbord},
  cell/.style={draw=mutec, fill=clifill, line width=0.3pt},
  ackcell/.style={draw=conbord, fill=confill, line width=0.3pt},
  guide/.style={draw=conbord, line width=0.3pt, dashed},
  marker/.style={font=\sffamily\scriptsize\bfseries, text=focusc},
  lbl/.style={font=\sffamily\footnotesize, text=inkc},
  bitlbl/.style={font=\sffamily\scriptsize, text=mutec}
]

% UM10204 Rev 7.0 3.1.4: START = SDA high->low while SCL high,
%                        STOP  = SDA low->high while SCL high.
% UM10204 Rev 7.0 3.1.5/3.1.6: each byte is 8 bits followed by an
% acknowledge bit on the ninth clock pulse.
% One bit slot = 0.35 wide. Slot j (j = 0..17) starts at x = 0.6 + j*0.35.
% Slots 0-7 address+R/W, slot 8 ACK, slots 9-16 data, slot 17 ACK.
% SCL is high in the middle of every slot; SDA only changes while SCL is low.
\def\ysda{2.0}
\def\yscl{0.0}
\def\w{0.35}

% ---------------- SDA ----------------
\node[lbl, anchor=east] at (-0.3, \ysda + 0.3) {SDA};
% idle high, START: high -> low while SCL is still high
\draw[sda] (-0.5, \ysda + 0.7) -- (0.2, \ysda + 0.7)
  -- (0.2, \ysda) -- (0.6, \ysda);
% address byte: 8 slots
\draw[cell] (0.6, \ysda) rectangle (3.4, \ysda + 0.6);
\node[bitlbl] at (2.0, \ysda + 0.3) {addr (7-bit + R/W)};
% ACK slot (9th clock)
\draw[ackcell] (3.4, \ysda) rectangle (3.75, \ysda + 0.6);
\node[bitlbl, text=conbord] at (3.575, \ysda + 0.3) {A};
% data byte: 8 slots
\draw[cell] (3.75, \ysda) rectangle (6.55, \ysda + 0.6);
\node[bitlbl] at (5.15, \ysda + 0.3) {data byte};
% ACK slot (9th clock)
\draw[ackcell] (6.55, \ysda) rectangle (6.9, \ysda + 0.6);
\node[bitlbl, text=conbord] at (6.725, \ysda + 0.3) {A};
% STOP: SDA low while SCL low, then low -> high after SCL has gone high
\draw[sda] (6.9, \ysda) -- (7.45, \ysda)
  -- (7.45, \ysda + 0.7) -- (8.0, \ysda + 0.7);

% ---------------- SCL ----------------
\node[lbl, anchor=east] at (-0.3, \yscl + 0.3) {SCL};
% idle high through START, then low before the first bit
\draw[scl] (-0.5, \yscl + 0.7) -- (0.45, \yscl + 0.7)
  -- (0.45, \yscl) -- (0.6, \yscl);
% 18 slots: low at the slot edges, one high pulse in the middle of each slot
\foreach \j in {0, ..., 17} {
  \pgfmathsetmacro{\xs}{0.6 + \j*\w}
  \pgfmathsetmacro{\xr}{\xs + 0.09}
  \pgfmathsetmacro{\xf}{\xs + 0.26}
  \pgfmathsetmacro{\xe}{\xs + \w}
  \draw[scl] (\xs, \yscl) -- (\xr, \yscl) -- (\xr, \yscl + 0.7)
    -- (\xf, \yscl + 0.7) -- (\xf, \yscl) -- (\xe, \yscl);
}
% after the last ACK: SCL low, then high before SDA rises (STOP)
\draw[scl] (6.9, \yscl) -- (7.15, \yscl)
  -- (7.15, \yscl + 0.7) -- (8.0, \yscl + 0.7);

% clock numbers under SCL: 1-8 per byte, 9 = ACK clock
\foreach \j/\n in {0/1, 1/2, 2/3, 3/4, 4/5, 5/6, 6/7, 7/8,
                   9/1, 10/2, 11/3, 12/4, 13/5, 14/6, 15/7, 16/8} {
  \node[bitlbl] at (0.6 + \j*\w + 0.175, \yscl - 0.3) {\n};
}
\node[bitlbl, text=conbord] at (3.575, \yscl - 0.3) {9};
\node[bitlbl, text=conbord] at (6.725, \yscl - 0.3) {9};

% ACK slot guides: SDA A cell <-> 9th SCL pulse
\draw[guide] (3.575, \ysda - 0.1) -- (3.575, \yscl + 0.8);
\draw[guide] (6.725, \ysda - 0.1) -- (6.725, \yscl + 0.8);

% START / STOP markers
\node[marker] at (0.2, \ysda + 1.0) {START};
\draw[->, draw=focusc, line width=0.4pt] (0.2, \ysda + 0.85) -- (0.2, \ysda + 0.75);
\node[marker] at (7.45, \ysda + 1.0) {STOP};
\draw[->, draw=focusc, line width=0.4pt] (7.45, \ysda + 0.85) -- (7.45, \ysda + 0.75);

\end{tikzpicture}
\end{document}
